\section{Cross-validation and checkpoint transport}
\label{sec:comparison}

Two independent checks constrain the method externally: a cross-validation
against PIConGPU's exact-quadrature field reference, and a checkpoint and
restart transport. Neither is a truth comparison.

\subsection{Cross-validation against the PIConGPU reference}

The reference is the TSC-convolved field (Section~\ref{sec:accuracy}),
evaluated in PIConGPU in its reference-only field-initializer mode
\citep{W7t}. picBunch is at
\texttt{339ca7ad} (\texttt{experimental/mvp}); \texttt{ctest} passes 10/10
\citep{W7r}. Both sides use identical Yee offsets, grid, and origin; the
reference staggered point maps to a picBunch node by index with zero snap
deviation \citep{W7r}.

At $318$ staggered points on a $64\times64\times32$ grid, cell
$10^{-7}$\,m, the overall worst per-component relative error is
$4.86\times10^{-2}$ for $\gamma=1$ ($E/z$) and $6.20\times10^{-1}$ for
$\gamma=10$ ($E/z$; $E/y$ is $5.16\times10^{-1}$), against a reference
reproducibility floor of $1.19\times10^{-6}$ ($E$) and
$1.43\times10^{-6}$ ($B$) \citep{W7r,D13}. picBunch is about four to six
orders of magnitude above that floor, a no-agreement statement at the
reference floor, not a defect: B is exactly zero for $\gamma=1$ in both, the
bulk agrees to below $1$--$3\%$, and the error is concentrated at the core
and the beamed peaks \citep{W7r,D13}.

The whole-field comparison, $N=1000$ on $32\times32\times16$, cell
$2\times10^{-7}$\,m, $16384$ nodes per component, gives the same picture:
$\gamma=1$ bulk median $3.8$--$7.8\times10^{-3}$ with an L2 of
$4.6$--$5.7\times10^{-2}$; $\gamma=10$ bulk median $1.4$--$3.3\times10^{-2}$
but a beamed-core overshoot that pushes the $E/y$, $E/z$ L2 to
$0.76$--$0.92$ \citep{W7r}. The forward $E/x$ is well resolved at
$\gamma=10$ ($\sim6\%$); the beamed transverse $E/y$, $E/z$ are
$50$--$130\%$ worst at the core, where the Lorentz-contracted field is a
sub-cell pancake that TSC and the spectral near-field kernel cannot resolve
\citep{W7r}.

\begin{table}[t]
\centering
\small
\caption{Whole-field comparison of picBunch against the PIConGPU
\texttt{referenceOnly} field, $N=1000$, $32\times32\times16$, cell
$2\times10^{-7}$\,m,
$16384$ nodes/component, rosi\_v100 \citep{W7r}. \texttt{worst\_rel} is
peak-normalized; median and p99 are per-point over nodes with
$|\mathrm{ref}|>10^{-3}\times$peak; L2 is pooled. Not a truth comparison.}
\label{tab:wholefield}
\begin{tabular}{lrrrr}
\toprule
component & worst\_rel & median & p99 & L2 \\
\midrule
\multicolumn{5}{l}{\emph{$\gamma=1$} (B identically 0)} \\
$E/x$ & $1.48\times10^{-1}$ & $3.92\times10^{-3}$ & $3.20\times10^{-1}$ & $4.62\times10^{-2}$ \\
$E/y$ & $1.49\times10^{-1}$ & $3.76\times10^{-3}$ & $2.93\times10^{-1}$ & $4.75\times10^{-2}$ \\
$E/z$ & $1.22\times10^{-1}$ & $7.81\times10^{-3}$ & $3.65\times10^{-1}$ & $5.67\times10^{-2}$ \\
\midrule
\multicolumn{5}{l}{\emph{$\gamma=10$} (B/x identically 0)} \\
$E/x$ & $1.36\times10^{-1}$ & $4.80\times10^{-3}$ & $1.23\times10^{-1}$ & $4.44\times10^{-2}$ \\
$E/y$ & $1.98$  & $1.36\times10^{-2}$ & $7.25$ & $7.63\times10^{-1}$ \\
$E/z$ & $1.93$  & $3.33\times10^{-2}$ & $10.1$ & $9.19\times10^{-1}$ \\
$B/y$ & $4.42\times10^{-1}$ & $3.12\times10^{-2}$ & $1.41$ & $2.58\times10^{-1}$ \\
$B/z$ & $4.63\times10^{-1}$ & $1.52\times10^{-2}$ & $0.93$ & $2.20\times10^{-1}$ \\
\bottomrule
\end{tabular}
\end{table}

\paragraph{Caveat: the sampling groups are distance-from-centroid.} The
$318$-point groups (inside-support, face, outside) were cut by distance from
the bunch centroid, in a bunch with a $\sim6$-cell spread \citep{D13}. That
grouping may understate how close the ``outside'' points actually are to a
source's TSC support, so the per-group numbers are a spatial profile, not a
clean inside/outside-of-support split. The per-source dependence in
Section~\ref{sec:accuracy} uses the explicit support union and is the
measurement that separates the per-source floor from any superposition
contribution \citep{W7t}.

\subsection{Checkpoint format and bitwise arrival}

The checkpoint writer emits a PIConGPU openPMD field checkpoint matching
PIConGPU's own plugin schema: root layout, openPMD version 5.0, $E$, $B$
$x/y/z$ at the exact Yee positions, grid spacing and offset, and the
\texttt{picongpu\_idProvider} mesh metadata \citep{W7B}. The default
precision is uniform \texttt{float32} for every field and particle record, so
a stock single-precision PIConGPU can load the file \citep{IMPL3}.
\texttt{v0.5.0} copies every species from the input bunch and carries a single
checkpoint-wide \texttt{picongpu\_idProvider} range \citep{INT1,FIX1,INT2};
the \texttt{charge} and \texttt{mass} records are written rather than omitted,
and the species is a single particle patch (Section~\ref{sec:interface}). The
remaining differences against a genuine PIConGPU checkpoint are non-blocking
metadata \citep{W7B}.

Field arrival is bitwise. A fresh \texttt{v0.5.0} checkpoint written from a
two-species input (16 electrons and 16 ions, \texttt{picongpu\_idProvider}
startId $0$, nextId $32$) restarts in PIConGPU with exit $0$ under both
\texttt{patch} modes (\texttt{bbox}: offset $(7,8,3)$, extent $(8,10,6)$;
\texttt{domain}: offset $(0,0,0)$, extent $(24,24,12)$), all six field
components bitwise-equal, $\max_\mathrm{abs\ diff} = 0.0$, electrons with ids
$0$ to $15$, first-step displacement exactly $\beta c\,dt$ with
$\beta = 0.9949874367$ \citep{INT2}. A validation-only build with the field
solver set to \texttt{None} loads the checkpoint and writes step 1; all six
components are bitwise-equal, max ULP $0$ \citep{W7B}.

\subsection{PML passthrough and the restart fixes}

On a $24\times24\times12$ deck the restart is exit $0$ with
\texttt{--periodic 1 1 1}, and the $16$ electrons load; initialization drops
from the $\sim116$\,s neuralBunch path to $1.8$\,s because the checkpoint
carries the field \citep{W7B}. The same investigation found the defects in
PIConGPU's restart path that the fixes PR addresses \citep{D15}:
\texttt{RestartFieldLoader} never reads \texttt{unitSI},
\texttt{gridSpacing} or \texttt{gridGlobalOffset}; the PML slabs are
unconditional; a successful restart suppresses the initial dump; and the
openPMD backend does not convert dataset dtypes. The fixes PR
(\texttt{psychocoderHPC-bot/picongpu} PR \#25, draft) adds restart metadata
validation, unitSI conversion, optional PML slabs, and an explicit dtype
contract, verified by real restarts including a factor-2 unitSI case
\citep{D15}.
